\section{The exact incidence-treewidth-two constant}\label{sec:width-two}
\begin{theorem}[One-sided cycle coloring]\label{thm:one-sided-coloring}
In a finite simple bipartite graph $G=(V,F;I)$ of treewidth at most two,
one can color the factor side $F$ with two colors so that every cycle
whose factor vertices all have one color contains an even number of
factor vertices. Consequently every zero-one matrix with this bipartite
support has a partition of its rows into two totally unimodular matrices.
\end{theorem}
The all-cycle assertion is stronger than excluding induced odd-factor
cycles (balancedness). The stronger statement will supply total
unimodularity for arbitrary integer cardinality slabs.

\begin{proof}
We prove a two-terminal series-parallel invariant. A network starts with
an edge; series composition identifies one terminal of each of two
otherwise disjoint networks, and parallel composition identifies both
terminal pairs. Intermediate networks are simple and bipartite subgraphs
of the eventual graph. A coloring is good if its monochromatic-factor
cycles have even factor count. A terminal path is $c$-monochromatic when
all its factor vertices, including factor terminals, have color $c$.
Its parity counts those factor vertices modulo two.

For each network there is a parity $\epsilon\in\{0,1\}$ with all the
following witnesses available. They are possibly different colorings.
For $V,V$ terminals and either color $c$, an \emph{active} good coloring
has at least one $c$-monochromatic terminal path, all such paths have
parity $\epsilon$, and no path is monochromatic in the other color.
If $\epsilon=0$, a \emph{blocking} good coloring with neither type of
monochromatic terminal path is also available. For mixed terminals,
with either prescribed color $c$ of the factor terminal, an active good
coloring has paths of color $c$, all of parity $\epsilon$; a blocking
witness with that prescribed terminal color is available whenever there
is no direct terminal edge. For $F,F$ terminals, with either prescribed
common color $c$, the same active witness exists; if $\epsilon=1$ a
blocking witness with those colors also exists. In addition, for either
assignment of distinct terminal colors, a good coloring exists. Such
colors automatically block monochromatic terminal paths. Finally a
network with a direct terminal edge is mixed and uses $\epsilon=1$.

The base edge satisfies the mixed invariant with parity one. For a
series composition let $\delta$ be one if the shared vertex is a factor,
and zero otherwise. The output parity is
$\epsilon=\epsilon_1\mathbin\oplus\epsilon_2\mathbin\oplus\delta$.
Every path concatenates component terminal paths; subtracting the
shared factor once accounts for $\delta$. Every cycle stays in a
component. Active witnesses concatenate by using the same desired color
and matching the shared factor's color when present. The required
additional witnesses have the following complete case analysis.
\begin{enumerate}
\item For outer types $V,V$ and a shared variable, use opposite active
colors in the two $V,V$ components. No monochromatic path can cross both.
\item For outer types $V,V$ and a shared factor, blocking is required
only if the component parities differ. The even mixed component has no
direct edge and can block with the prescribed shared-factor color;
use an active witness of that color in the other component.
\item For mixed outer types and a shared variable, use the active color
opposite to the outer factor's color in the $V,V$ component and use
that prescribed factor color in the mixed component. This blocks.
\item For mixed outer types and a shared factor, give the shared factor
the color opposite to the outer factor. The $F,F$ component uses its
distinct-terminal witness, and the mixed component uses a compatible
active witness. This again blocks.
\item For outer types $F,F$ of common color and a shared variable,
blocking is required only for output parity one, so the two mixed
component parities differ. Block the even component, which has no
direct edge, and activate the other, respecting both outer colors.
\item For outer types $F,F$ of common color and a shared factor, color
the shared factor oppositely. Both $F,F$ components have their required
distinct-terminal good witnesses and block.
\item For outer $F,F$ terminals of distinct colors and a shared
variable, activate each mixed component in its outer terminal's color.
If instead the shared vertex is a factor, give it either color and use
in each component the equal-color active or distinct-color witness
as appropriate.
\end{enumerate}
These cases include reversed mixed types and exhaust all triples of
terminal types. A nontrivial series composition has no direct outer edge,
so its mixed blocking requirement has also been checked.

In a parallel composition every new simple cycle consists of one
terminal path from each component. If the paths are monochromatic in
one color, its factor parity is their parity XOR plus $\tau$, the
number of factor terminals modulo two. Thus $\tau=0$ for equal terminal
types and $\tau=1$ for mixed types. We check the invariant as follows.
For $V,V$ terminals with equal component parities, activate both in the
same color. If both are even, both may also block. For different parities,
activate the odd component and block the even component. Output parity
is $\epsilon_1\lor\epsilon_2$; whenever both components contribute paths
their equal parities make new cycles even.
For $F,F$ terminals of common prescribed color, activate both when
parities agree; if both are odd, both may also block. If parities differ,
activate the even component and block the odd one. Output parity is
$\epsilon_1\land\epsilon_2$. With distinct prescribed terminal colors,
use their good witnesses in both components; a crossing monochromatic
cycle would contain both differently colored terminals and is impossible.
For mixed terminals with no direct edge, both components can block in
the prescribed factor color. Activate either one and block the other
for an active witness, choosing its parity as output, or block both
for the blocking witness. For mixed terminals with a direct edge,
simplicity implies exactly one component contains it. Activate its
parity-one witness and block the edge-free component. No blocking
output is required. These mixed choices create no crossing monochromatic
cycle. The induction is complete.

A treewidth-two graph has no $K_4$ minor, hence no subdivision of $K_4$.
The classical block characterization supplies a two-terminal
series-parallel representation for each nontrivial block
\citep[Theorem 3.1]{HassinTamir1986}. Apply the invariant to each block.
Root the block-cut forest. If a child block attaches at a factor,
interchange its two colors if needed to match the articulation color;
if it attaches at a variable no matching is needed. Bridges use the
base edge and isolated factors receive arbitrary colors. Every cycle
lies in one block, proving the graph assertion.

For the matrix assertion use Camion's criterion: a $0,\pm1$ matrix
is totally unimodular exactly when every Eulerian submatrix has entry
sum divisible by four \citep[Theorem 6.5]{Cornuejols2000}. In either
color class an Eulerian submatrix has even row and column degrees.
Its bipartite support decomposes into edge-disjoint simple cycles.
Each has an even number of factor vertices, hence length divisible
by four. Its total number of ones is consequently divisible by four,
as required. The argument covers rectangular Eulerian submatrices too.
\end{proof}

\paragraph{Formal proof of the matrix step.}
The \href{../formal/topics/19-structural-multilinear/COVERAGE.md}{topic 19 formalization}
uses a Ghouila--Houri signing proof in place of Camion's criterion.
For each selected set of rows, pair its entries within each column, leaving
at most one unpaired row. A cycle in the graph of paired rows lifts to an
incidence trail. The all-cycle parity property forces every such cycle to
be even, so the pair graph admits two colors. Opposite signs on each pair
leave every column sum in $\{-1,0,1\}$. The sufficiency of this signing
condition for total unimodularity is proved in Lean. The formal development
also constructs the factor coloring from a tree decomposition of width two
and proves the cardinality-slab integrality used below; it does not assume
these intermediate conclusions. The
\href{../formal/topics/19-structural-multilinear/VERIFICATION.md}{verification record}
records the checks and source snapshots.

\begin{theorem}[Sharp width-two gap]\label{thm:width-two-gap}
A sum of discrete-convex cardinality factors with incidence treewidth
at most two satisfies $\tgap\le2\hgap$. This includes positive
monomials on unit cubes and zero-lower boxes, and original positive
products with a common aspect ratio within each scope. The sharp
supremum is two, including for every fixed common positive aspect ratio
$\rho>1$.
\end{theorem}
\begin{proof}
Partition factor rows into two TU classes by
Theorem~\ref{thm:one-sided-coloring}. For a class $C$ and prescribed
means $p$, the polytope
\[
 0\le z\le1,\qquad
 \Bigl\lfloor\sum_{i\in e}p_i\Bigr\rfloor
 \le\sum_{i\in e}z_i\le
 \Bigl\lceil\sum_{i\in e}p_i\Bigr\rceil\quad(e\in C)
\]
is integral and contains $p$. Indeed row sign changes, duplicate rows,
and unit rows preserve total unimodularity; each basic solution with
integer right-hand side is integral by its determinant $\pm1$.
A binary-vertex decomposition of $p$ attains all class lower envelopes
by Lemma~\ref{lem:cardinality-envelopes}. Under that law every other
factor has nonnegative expected deficiency from its own upper value.
The equal mixture of the two class laws therefore attains at least
half the sum of local gaps. Their common upper extremizer makes this
a lower bound for half the full termwise gap in the scalar hull gap.

For positive unit-cube monomials one can equivalently use failure rows
$\sum_{i\in e}z_i\le1$ when their mean sum is at most one and
$\sum_{i\in e}z_i\ge1$ otherwise. A classwise integral decomposition
attains the failure union bound in every row. The mixed unit-right-hand-side
integrality theorem for balanced matrices suffices for this particular
argument \citep[Theorem 6.13]{Cornuejols2000}; it would not justify
arbitrary integer cardinality slabs without the stronger TU conclusion.
Common-aspect products are cardinality factors by
Corollary~\ref{cor:cardinality-box}. Proposition~\ref{prop:feedback-sharp}
proves the unit-cube lower supremum two.

For sharpness at a fixed positive ratio, put $\epsilon=1/\rho$ and
$\alpha=1-\epsilon$. On $[\epsilon,1]^{n+1}$ use
\[
 f_n(a,x)=\alpha^{-1}a\sum_{i=1}^nx_i+\prod_{i=1}^nx_i,
 \qquad \E A=1/n,\quad\E X_i=1-1/n,
\]
where $a=\epsilon+\alpha A$ and $x_i=\epsilon+\alpha X_i$ are endpoint
variables. The normalized nonlinear bilinear terms are $\alpha AX_i$,
so their gaps sum to $\alpha$. The large product is $\epsilon^R$ for
failure count $R$, with $\E R=1$. Its lower value is $\epsilon$ by
adjacent-cardinality attainment, and its upper value
$1-(1-\epsilon^n)/n$. Hence
\[
 \tgap_n=2\alpha-\frac{1-\epsilon^n}{n},\qquad
 \hgap_n=\max\E[\alpha AR+1-\epsilon^R]-\frac{1-\epsilon^n}{n}.
\]
For any $M>0$, $1-\epsilon^r\le\alpha r$ on integer $r\ge0$ gives
\[
 \E[1-\epsilon^R]\le\alpha(1-\E[R;R>M])+1/M,
 \quad \alpha\E[AR]\le\alpha M/n+\alpha\E[R;R>M].
\]
Markov's inequality was used in the first bound. Thus
$\hgap_n\le\alpha+\alpha M/n+1/M$. Exactly one uniformly failed leaf
and an independent anchor give
$\hgap_n\ge\alpha+\alpha/n-(1-\epsilon^n)/n$.
Taking $M=\sqrt n$ yields $\hgap_n\to\alpha$ and
$\tgap_n/\hgap_n\to2$. The incidence graph is the same width-two graph
as before; all means are interior. Scaling physical coordinates by
$\rho$ maps the box to $[1,\rho]^{n+1}$ and preserves the gaps of the
correspondingly rescaled positive polynomial.
\end{proof}

\subsection{What fails for general nonnegative factors}
\begin{example}[An exact width ratio of three]\label{ex:parity-three}
Use six fair variables in pairs $AB,BC,CA$, and put
\[
 E(u,v)=1-u-v+2uv,\qquad N(u,v)=u+v-2uv.
\]
These are the multiaffine equality and inequality indicators, both in
$[0,1]$ on the square. Define
\[
 g_A=E(AB)E(CA),\qquad g_B=N(AB)E(BC),\qquad g_C=N(BC)N(CA).
\]
Each local upper envelope is one at the center: mix a satisfying
assignment and its complement. Each local lower envelope is zero,
by the analogous mixture of a falsifying assignment. No two factors
can equal one on a binary assignment, because adjacent factors demand
opposite parities on their common pair. Thus the sum's upper envelope
is at most one, attained by the all-zero/all-one mixture, which
satisfies $g_A$. The pair assignment $AB=01,BC=01,CA=00$ and its
complement make all factors zero, proving the sum's lower envelope
zero. Therefore $\tgap=3$ and $\hgap=1$ for these three entire factors.

A central bag $\{A,B,C\}$ and six attached bags
$\{A,B,X_{AB,j}\}$, $\{B,C,X_{BC,j}\}$, $\{C,A,X_{CA,j}\}$,
$j=1,2$, form a width-two decomposition. A four-cycle proves the width
is exactly two. Every variable has frequency two. The factors are
nonnegative on the full cube, but their monomial coefficients have mixed
signs, so neither the positive-monomial nor the convex-cardinality
factor theorem applies.
\end{example}

\begin{proposition}[Independent-set payoff family]\label{prop:independent-payoffs}
For a simple graph $G$ with nonnegative vertex weights $w_v$, place two
fair bits on each edge and assign its endpoints opposite demanded
parities. Give vertex $v$ payoff $w_v$ if all its incident demanded
parities hold, and zero otherwise. The sum of separate \emph{maximum
expected payoffs} is $\sum_vw_v$, while the global maximum is
$\alpha_w(G)$, the maximum independent-set weight. With all constructed
factor scopes retained, including zero-weight factors, a graph with
an edge has incidence treewidth $\max\{2,\tw(G)\}$.
\end{proposition}
\begin{proof}
Each local satisfaction pattern and its complement attain $w_v$ with
fair means. Satisfied vertices form an independent set, since the two
endpoints of an edge demand opposite parities. Conversely fix a maximum
independent set, give each incident edge its selected endpoint's demand,
and choose remaining parities arbitrarily. Complement-pair mixing
satisfies all selected factors with fair means, proving the maximum
$\alpha_w(G)$; any extra satisfied factors still form an independent
set and cannot improve that maximum.
To a decomposition of $G$ attach two size-three bags $\{u,v,X_{uv,j}\}$
for every edge at a bag containing $u,v$. This proves the treewidth
upper bound. Contracting one length-two path per edge and deleting
the other gives $G$ as a minor; each edge also creates a four-cycle,
proving both lower bounds.
\end{proof}
Thus $K_{k+1}$ with unit vertex weights gives maximum-payoff ratio $k+1$
at incidence width $k$ for $k\ge2$. A general payoff maximum ratio is not an envelope-width
ratio: the minima must also be computed. Example~\ref{ex:parity-three}
does compute them and gives the claimed exact width obstruction.

\subsection{Limits of the coloring method}\label{subsec:width-three-boundary}
A distinct open question at incidence treewidth two asks whether the
factors admit two classes whose incidence subgraphs are each chordal
bipartite, meaning they contain no induced cycle of length at least six.
The proved all-cycle parity property permits induced eight-cycles and
does not establish this target. The recorded experiment found such a
partition for all 5,000 generated supports; this is finite evidence only.

The all-cycle coloring invariant cannot extend with a bounded number
of colors to width three. In $K_{3,m}$, $m\ge3$, with its three-vertex
side interpreted as variables, every three factors lie on a six-cycle.
Each permitted color class therefore contains at most two factors,
requiring $\lceil m/2\rceil$ colors, although the graph has treewidth
three. Its all-ones matrix is already TU: determinants of order at
least two vanish. The six-cycles have chords and do not obstruct
balancedness. This separates the stronger coloring mechanism from a
TU-row-partition target.

Finite searches found three balanced classes in 3,000 generated
width-at-most-three supports (212 required more than two), and three
TU classes in 300 further supports. Their accepted widths were
certified by particular elimination orders; rejected greedy orders did
not certify larger width. The TU tests enumerated column signings by
the Ghouila-Houri criterion and used exact coloring backtracking.
These searches are finite evidence, not proofs of universal partitions.
The exact $W_k$ for $k\ge3$, $k$-class balanced or TU partitions, and
sharp planar constants remain open here. A planar bipartite incidence
graph does admit an orientation of maximum outdegree at most two:
Hall's theorem assigns every edge to one of two slots at an endpoint,
since every edge set uses at most twice as many edges as incident
vertices. Orient from the assigned endpoint. Thus the general
orientation theorem gives a finite planar bound, without settling its
sharp value. Auxiliary reductions and the simpler forest-coloring
obstruction are recorded in Appendix~\ref{app:structural-auxiliary}.
